% typeset: Pdftex
% Afterwards compile with pdflatex > bibtex > pdflatex > pdflatex.
% in TeXShop preferences, change edit from bibtex to biber
% https://tex.stackexchange.com/questions/144407/no-items-shown-bibliography-beamer?rq=1

% https://tex.stackexchange.com/questions/1423/is-there-a-nice-way-to-compile-a-beamer-presentation-without-the-pauses
%\documentclass[handout]{beamer}

%! LaTeX Error: File `/Volumes/repos/bitbucket/strange/shared/collection/lib-colors.tex' not found.
% edit 
% nursery-slide-decks/beamer/dtra/global/setup-global

\documentclass[]{beamer}
% point to content libraries
\newcommand{\pGlobal}			{../../../global/}
\newcommand{\pGlobalSetup}		{\pGlobal setup-global/}	

% load LaTeX packages and personal macros
\input{\pGlobalSetup setup-global-beamer}
\input{\pLocalSetup setup-local}

\setbeamertemplate{bibliography item}{\insertbiblabel}
\usepackage[style = numeric]{biblatex}
%\usepackage[square,numbers]{natbib}
%\setcitestyle{numbers}
%\bibliography{\pSections em-gauge.bib}
%\usepackage{auto-pst-pdf}
%\usepackage{pst-solides3d}
\bibliography{\pSections redux.bib}

\title
[Taylor's Fireball Analysis]
{Taylor's Fireball Analysis\\Reproducing Numerical Results}

\author[\dtriacAlpha NIAC Tech Staff]{\Topa \\ \TopaEmail}
\institute{\dtriac} 
\medskip

\date{\today}

\begin{document}

\begin{frame}
	\titlepage
	%\distA
\end{frame}
%
\begin{frame}\frametitle{Search for Similarity}
\begin{table}[htp]
%\caption{default}
\begin{center}
	\begin{tabular}{cc}
	\href{http://homepage.physics.uiowa.edu/~ghowes/teach/phys7729/docs/Taylor50b.pdf\#page=13}{
	\begin{overpic}[ scale = 0.25 ]
		{\pLocalGraphics taylor/taylor-image}
		%\put(30, 10)	{\colorbox{white}{ $R = S(\gamma) t^{\tfrac{2}{5}} E^{\tfrac{1}{5}} \rho_{0}^{-\tfrac{1}{5}} $ }}
		%\put(33, 30)	{\colorbox{white!10}{ $R = S(\gamma) \sqrt[5]{\frac{E t^{2}}{\rho} } $ }}
		\put(43, 4)		{ \color{white}{$100 $ \textrm{m}} }
		\put(33, 30)	{ \bl{$R = S(\gamma) \sqrt[5]{\frac{E t^{2}}{\rho_{0}} } $} }
	\end{overpic}}
	\end{tabular}
\end{center}
%\label{tab:label}
\end{table}%
\end{frame}

%
\begin{frame}\frametitle{Photometric Yield Determination}
\begin{table}[htp]
%\caption{default}
\begin{center}
	\begin{tabular}{cc}
	\href{http://homepage.physics.uiowa.edu/~ghowes/teach/phys7729/docs/Taylor50b.pdf\#page=13}{
	\begin{overpic}[ scale = 0.25 ]
		{\pLocalGraphics taylor/taylor-image}
		%\put(30, 10)	{\colorbox{white}{ $R = S(\gamma) t^{\tfrac{2}{5}} E^{\tfrac{1}{5}} \rho_{0}^{-\tfrac{1}{5}} $ }}
		%\put(33, 30)	{\colorbox{white!10}{ $R = S(\gamma) \sqrt[5]{\frac{E t^{2}}{\rho} } $ }}
		\put(43, 4)		{ \color{white}{$100 $ \textrm{m}} }
		\put(33, 30)	{ \bl{$R = S(\gamma) \sqrt[5]{\frac{E t^{2}}{\rho_{0}} } $} }
	\end{overpic}}
	\end{tabular}
\end{center}
%\label{tab:label}
\end{table}%
\end{frame}

\begin{frame}\frametitle{Overview}
	\tableofcontents[hideallsubsections]
\end{frame}
%
%   %   %   %   %   %   %   %   %   %   %   %   %   %   %
\section{Photometric Yield Determination}
\begin{frame}\frametitle{Trinity Test}
\begin{table}[htp]
%\caption{default}
\begin{center}
	\begin{tabular}{cc}
	\href{http://homepage.physics.uiowa.edu/~ghowes/teach/phys7729/docs/Taylor50b.pdf\#page=13}{
	\begin{overpic}[ scale = 0.25 ]
		{\pLocalGraphics taylor/taylor-image}
		%\put(30, 10)	{\colorbox{white}{ $R = S(\gamma) t^{\tfrac{2}{5}} E^{\tfrac{1}{5}} \rho_{0}^{-\tfrac{1}{5}} $ }}
		%\put(33, 30)	{\colorbox{white!10}{ $R = S(\gamma) \sqrt[5]{\frac{E t^{2}}{\rho} } $ }}
		\put(43, 4)		{ \color{white}{$100 $ \textrm{m}} }
		\put(33, 30)	{ \bl{$R = S(\gamma) \sqrt[5]{\frac{\rd{E} t^{2}}{\rho_{0}} } $} }
	\end{overpic}}
	\end{tabular}
\end{center}
\label{tab:R}
\end{table}%
\end{frame}

%
\begin{frame}\frametitle{Trinity Test}
\begin{table}[htp]
%\caption{default}
\begin{center}
	\begin{tabular}{cc}
	\href{http://homepage.physics.uiowa.edu/~ghowes/teach/phys7729/docs/Taylor50b.pdf\#page=13}{
	\begin{overpic}[ scale = 0.25 ]
		{\pLocalGraphics taylor/taylor-image}
		%\put(30, 10)	{\colorbox{white}{ $R = S(\gamma) t^{\tfrac{2}{5}} E^{\tfrac{1}{5}} \rho_{0}^{-\tfrac{1}{5}} $ }}
		%\put(33, 30)	{\colorbox{white!10}{ $R = S(\gamma) \sqrt[5]{\frac{E t^{2}}{\rho} } $ }}
		\put(43, 4)		{ \color{white}{$100 $ \textrm{m}} }
		\put(33, 30)	{ \bl{$\rd{E} = \bl{\rho_{0} \cdot L \cdot R^{5}t^{-2}}$} }
	\end{overpic}}
	\end{tabular}
\end{center}
\label{tab:E}
\end{table}%
\end{frame}
%


%   %   %   %   %   %   %   %   %   %   %   %   %   %   %
%\section{Scattered Results}
\begin{frame}\frametitle{Taylor's Method: Yield in Kilotons}
\center
	\href{https://www.osti.gov/opennet/manhattan-project-history/publications/DOENuclearTests.pdf\#page=25}{
	\begin{overpic}[ scale = 1.0 ]
		{\pLocalGraphics niac/yield-line-fine}
		\put(33, 20)	{\colorbox{white}{$E = \rho_{0} \cdot  L \cdot R^{5}t^{-2}$}}
	\end{overpic}}
	%\ \\
	%\fullcite{united2000united}
\end{frame}


\begin{frame}\frametitle{Numeric Details}
\center
	%\href{}{
	\begin{overpic}[ scale = 0.45 ]
		{\pLocalGraphics niac/topa-tbl-01}
		%\put(50, 50)	{\colorbox{white}{$a+b$}}
	\end{overpic}
\end{frame}


\begin{frame}\frametitle{Taylor Solution}
\center
	%\href{}{
	\begin{overpic}[ scale = 0.5 ]
		{\pLocalGraphics niac/theorem}
		%\put(50, 50)	{\colorbox{white}{$a+b$}}
	\end{overpic}
\end{frame}

\begin{frame}\frametitle{Proof\jumpBig}
\center
	%\href{}{
	\begin{overpic}[ scale = 0.5 ]
		{\pLocalGraphics niac/proof}
		%\put(50, 50)	{\colorbox{white}{$a+b$}}
	\end{overpic}
\end{frame}

%   %   %   %   %   %   %   %   %   %   %   %   %   %   %
\section{Power Law}
%
\begin{frame}\frametitle{Power Law is Nonlinear}
		%
\begin{equation}
	R = S(\gamma) \sqrt[5]{\frac{E t^{2}}{\rho_{0}} }
\label{eq:Taylor}
\end{equation}
		%
	$$ \Downarrow $$
		%
\begin{equation}
	R(t) = a t^{\tfrac{2}{5}}
\label{eq:power-law}
\end{equation}
		%
		%
	$$ \Downarrow $$
		%
\begin{equation}
	R(t) = a t^{b}
\label{eq:power-law-b}
\end{equation}
		%
\end{frame}
%
%
\begin{frame}\frametitle{Solution Curves}
\center
	%\href{}{
	\begin{overpic}[ scale = 1.0 ]
		{\pLocalGraphics mathematica/Taylor-redux-true}
		\put(33, 50)	{\colorbox{white}{Taylor solution}}
		\put(60, 40)	{\colorbox{white}{nonlinear solution}}
	\end{overpic}
\end{frame}
%
\begin{frame}\frametitle{Comparing Solutions}
\center
	%\href{}{
	\begin{overpic}[ scale = 0.65 ]
		{\pLocalGraphics mathematica/nonlinear-contour-zoom}
		\put(7, 87)	{\colorbox{white}{$\normt{at^{b} - R}$}}
		\put(49, -5)			{\colorbox{white}{$a$}}
		\put(-7, 48)			{\colorbox{white}{$b$}}
	\end{overpic}
\end{frame}
%
%\begin{frame}\frametitle{Mathematica Solution}
%\center
%	%\href{}{
%	\begin{overpic}[ scale = 0.4 ]
%		{\pLocalGraphics mathematica/mm-screen-02}
%		%\put(50, 50)	{\colorbox{white}{$a+b$}}
%	\end{overpic}
%\end{frame}
%
\begin{frame}\frametitle{Solution Using \href{https://www.wolfram.com/mathematica/	}{Mathematica}}
\center
	\href{run:../local/articles-redux/2023-film-working-group.pdf}{
	\begin{overpic}[ scale = 0.375 ]
		{\pLocalGraphics mathematica/norm-min}
		\put(55, 5)	{\colorbox{white}{$(\bl{a}, \bl{b})_{LS}\equiv \argmin\limits_{(\bl{a},\bl{b})\in\mathbb{R}^{2}} \normt{\bl{a} t^{\bl{b}} - R}$}}
	\end{overpic}}
\end{frame}
%
\begin{frame}\frametitle{Individual Points or Aggregation\jumpBig}
\begin{figure}[htbp]
\begin{center}
		\vspace{0.45cm}
		\begin{overpic}[ scale = 0.55 ]
		{\pLocalGraphics mathematica/redux-yield}
		\put(8, 102)	{\colorbox{white}{$E_{LS} = \overline{E} = 16.7 \pm 3.2 Kt$}}
		\put(-12, 34)	{\rotatebox{90}{\colorbox{white}{Yield, Kt}}}
		\put(35, -9)	{\colorbox{white}{$\log_{10} t$}}
		\end{overpic}
		\vspace{0.35cm}
%\caption{Point measurements and yield using the updated parameters. The thick line represents the least squares value of $\cxi = 6.42205\expo{23}$ erg.}
\label{fig:fit-yield}
\end{center}
\end{figure}
\end{frame}
%
%   %   %   %   %   %   %   %   %   %   %   %   %   %   %
\section{Stack Overflow}
\begin{frame}\frametitle{\href{https://math.stackexchange.com/questions/670654/least-squares-solution-for-y-axb-after-logarithmic-transformation/2241027\#2241027}{
Full Example Online}}
\center
	%\href{https://math.stackexchange.com/questions/670654/least-squares-solution-for-y-axb-after-logarithmic-transformation/2241027\#2241027}{
	\href{run:../local/articles/Mathematics-Stack-Exchange.pdf}{
	\begin{overpic}[ scale = 0.325 ]
		{\pLocalGraphics mse}
	\end{overpic}}
\end{frame}
%

%
\begin{frame}\frametitle{Further Examples\jumpBig}
\begin{itemize}
	\item \href{https://math.stackexchange.com/questions/1488747/least-square-approximation-for-exponential-functions/2163333\#2163333}{Least Square Approximation for Exponential Functions}
	\item \href{https://math.stackexchange.com/questions/1934347/substitution-yielding-linear-function-for-least-squares-fitting/2170544\#2170544}{Substitution yielding linear function for least squares fitting}
	\item \href{https://math.stackexchange.com/questions/699071/using-least-squares-method-to-determine-unknown-variables/2204288\#2204288}{Using Least Squares method to determine unknown variables}
\end{itemize}
\end{frame}
%\begin{frame}\frametitle{Similarity Implies \href{https://mathworld.wolfram.com/Power.html}{Power}  Law Model}
%\huge{
%\begin{equation}
%	\bl{R} = \bl{S(\gamma) \sqrt[5]{\frac{\rd{E}\, t^{2}}{\rho_{0}} }}
%\tag{\ref{eq:one}}
%\end{equation}
%}
%\end{frame}
%%
%\begin{frame}\frametitle{Similarity Implies \href{https://mathworld.wolfram.com/Power.html}{Power} Law Model}
%\huge{
%\begin{equation*}
%	R = \mg{S(\gamma)} \sqrt[5]{\frac{\mg{E}\, t^{2}}{\mg{\rho_{0}}} } \quad \Rightarrow \quad
%	R = \bl{a} t^{\bl{b}} 
%\end{equation*}
%}
%\end{frame}
%%
%\begin{frame}\frametitle{Similarity Implies \href{https://mathworld.wolfram.com/Power.html}{\yl{Power}}  \yl{Law} Model}
%% Taylor's premise
%\huge{
%\begin{equation}
%	R(\bl{a}, \bl{b},t) = \bl{a} t^\bl{b}
%\label{eq:model}
%\end{equation}
%}
%\end{frame}
%%
%\begin{frame}\frametitle{Least Squares \rd{Definition}}
%\huge{
%\begin{equation}
%	r(a, b) = \bl{a} t^{\bl{b}} - R
%\label{eq:residual}
%\end{equation}
%\pause
%\begin{equation}
%	(\bl{a}, \bl{b})_{LS} \rd{\, \equiv \,} \argmin_{(\bl{a},\bl{b})\in\mathbb{R}^{2}} \normt{\bl{a} t^{\bl{b}} - R}
%\label{eq:merit}
%\end{equation}
%}
%\end{frame}
%%
%\begin{frame}\frametitle{Solution Using \href{https://www.wolfram.com/mathematica/	}{Mathematica}}
%\center
%	%\href{}{
%	\begin{overpic}[ scale = 0.375 ]
%		{\pLocalGraphics Taylor/mm-screen-02}
%		%\put(35, 52)	{\colorbox{white}{\rd{Faux Linearization}}}
%	\end{overpic}
%\end{frame}
%%
%\begin{frame}\frametitle{Comparing Solution Points}
%\center
%	%\href{}{
%	\begin{overpic}[ scale = 0.6 ]
%		{\pLocalGraphics Taylor/Taylor-redux-both}
%		\put(7, 87)		{\colorbox{white}{$\normt{r(a,b)} = \normt{a t^{b} - R}$}}
%		\put(-9, 46)	{\rotatebox{90}{\colorbox{white}{$b$}}}
%		\put(49, -5)	{\colorbox{white}{$a$}}
%	\end{overpic}
%\end{frame}
%%
%\begin{frame}\frametitle{Comparing Solution Points}
%\center
%	%\href{}{
%	\begin{overpic}[ scale = 0.6 ]
%		{\pLocalGraphics Taylor/Taylor-redux-both}
%		\put(7, 87)		{\colorbox{white}{$r(a,b) = \normt{a t^{b} - R}$}}
%		\put(-9, 46)	{\rotatebox{90}{\colorbox{white}{$b$}}}
%		\put(49, -5)	{\colorbox{white}{$a$}}
%		\put(100, 47)	{\colorbox{white}{\bl{$r = 510$}}}
%		\put(100, 88)	{\colorbox{white}{\rd{$r =1456$}}}
%	\end{overpic}
%\end{frame}
% main sections
%	\input{\pSections "sec-background"}
%	\input{\pSections "sec-linear-least-squares"}
%	\input{\pSections "sec-nonlinear-least-squares"}
%	%\input{\pSections "sec-redux"}
%	\input{\pSections "sec-backup"}

%   %   %   %   %   %   %   %   %   %   %   %   %   %   %
\section{Bibliography}
%\subsection{References}
\begin{frame}[t,allowframebreaks]
\frametitle{Bibliography}
\begin{enumerate}
	\item \fullcite{taylor1950formationI}
	\item \fullcite{taylor1950formationII}
	\item \fullcite{sedov2018similarity}
%\end{enumerate}
%\end{frame}
%
%\begin{frame}[t,allowframebreaks]
%\frametitle{Bibliography}
%\begin{enumerate}
	\item \fullcite{deakin2011gi}
	\item \fullcite{diaz2021explosion}
	\item \fullcite{aouad2021beirut}
\end{enumerate}
\end{frame}
%\nocite{*}
%\printbibliography[heading=none]

\end{document}

%\tiny
%\scriptsize
%\footnotesize
%\small
%\normalsize
%\large
%\Large
%\LARGE
%\huge
%\Huge

%\, thin space (normally 1/6 of a quad);
%\> medium space (normally 2/9 of a quad);
%\; thick space (normally 5/18 of a quad);
%\! negative thin space (normally 1/6 of a quad).

\begin{frame}\frametitle{Frame Title}
\begin{enumerate}
	\item 
	\item 
	\item 
\end{enumerate}
\end{frame}

\begin{frame}\frametitle{Frame Title}
\center
	\href{url}{
	\begin{overpic}[ scale = 1.0 ]
	{\pLocalGraphics graphic-file}
		%\put(-7,-10){Auxiliary text.}
	\end{overpic}}
\end{frame}

\begin{frame}\frametitle{Frame Title}
\begin{table}[htp]
%\caption{default}
\begin{center}
	\begin{tabular}{cc}
		 %
		 %
		 %
	\end{tabular}
\end{center}
%\label{tab:label}
\end{table}%
\end{frame}
